\newpage

\section{toyRISC Generic Implementation}\label{toyRiscDesign}

The toyRISC processor described in section \ref{toyRISCprocessor} in implemented, simulated and tested in this section. The project, starting from the instruction set architecture described in file 0\_DEFINES.vh  (see \ref{lst:toyRISCdefines}), consists in:
\begin{itemize}
    \item RTL description in System Verilog
    \item a binary code generator, a simple assembler, used to generate the codes defined in Table \ref{tab:riscISA}
    \item a simulator module used to generate the environment for the processor under test
    \item a small set of simple programs used to validate the main mechanisms and instructions.
\end{itemize}

The structure of the toyRISC processor is a three-level hierarchy of circuits:
\begin{description}
    \item[toyRISC]: the very simple RISC processor used to illustrate the principle of a processing unit as a system based on a loop connected two automata (Program Counter and Register with ALU) serially connected to the module DCD containing an automaton, used to manage the interrupt signal, and few combinational circuits used as a small decode unit
        \begin{description}
            \item[DCDtoyRISC]: the complex part of toyRISC which is a very small 2-state automaton and few gates; has a quantitatively insignificant share in the structural economy of the entire system
            \item[PCtoyRISC]: is a simple circuit used to manage the program control
            \item[RALUtoyRISC]: is a simple circuit used as execution unit; from a quantitative point of view it has the highest share in the economy of the system
                \begin{description}
                    \item[regFile]: defined in \ref{registerFile} and generically described  in \ref{lst:A09_regFile.sv}
                    \item[ALU]: defined in \ref{alu1} and generically described  in \ref{lst:A08_alu.sv}
                    \item[bigMux]: defined in \ref{genericMux}
                \end{description}     
        \end{description}
\end{description}

\subsection{toyRISC: the top level}


The top level connections of our simple processor refers to:
\begin{itemize}
    \item synchronizing signals: {\tt clock} and {\tt reset}
    \item interrupt signals: {\tt int} and {\tt inta}
    \item program memory: address ({\tt nextPC} is limited to a 10-bit address to keep the simulation fast) and instruction ({\tt instr}); no read signal needed because the read operation is implicit in each clock cycle, and write is not considered in the simple simulation environment we consider, where the program is generated by the simulation environment using the {\tt RISCcodeGenerator.sv} inserted in {\tt testRISC} module.
    \item data memory: with 32-bit data in-out connection to a 1K  word memory (for fast simulation reasons), and the command bits.
\end{itemize}

\includeverilogcode[caption={toyRISC. File: toyRISC.sv.  (SystemVerilog)}, 
                    label={lst:toyRISC}, 
                    firstline=1, 
                    firstnumber=1, 
                    lastline=68]
                    {\verilogSourceDir/A17_toyRISC.sv}

\placedrawing[H]{figures/A17_RISCstructure.lp}{Top module of toyRISC processor.}{toyRiscProc}

In Figure \ref{toyRiscProc} is represented the schematic provided by the Vivado tool for the top module of toyRISC.


\subsubsection{Decode module: the small complex part} is the only complex part of torRISC. It contains a finite automaton and some glue gates Therefore, it is a complex circuit, but quantitatively it is insignificant.

\includeverilogcode[caption={Decoder DCDtoyRISC. File: DCDtoyRISC.sv.  (SystemVerilog)}, 
                    label={lst:DCDtoyRISC}, 
                    firstline=1, 
                    firstnumber=1, 
                    lastline=68]
                    {\verilogSourceDir/A17_DCDtoyRISC.sv}

\subsubsection{Program counter: the simple control automaton} circuit was defined in \ref{progCounter} and  a generic structure, designed in \ref{lst:A10_programCounter.sv}. In Listing \ref{lst:DCDtoyRISC} is described the behavior of the simple automaton responsible for the program address management. It is simple because the loop over the register consists only simple, recursive defined, circuits: adder, incrementer, multiplexer.

\includeverilogcode[caption={Program counter PCtoyRISC. File: PCtoyRISC.sv.  (SystemVerilog)}, 
                    label={lst:DCDtoyRISC}, 
                    firstline=1, 
                    firstnumber=1, 
                    lastline=68]
                    {\verilogSourceDir/A17_PCtoyRISC.sv}

\subsubsection{RALU: the big sized simple part} is the other simple automaton, with the maximum weight in processor size, defined in \ref{rALU} (see also Figure \ref{ralu} with a generic definition in Listing \ref{lst:A10_RALU.sv}, based on the functions listed in Listing \ref{lst:A10_defineRALU.vh} and on circuits described in the previous two chapters: ALU and register file. The 2-input multiplexer, from  the generic definition, is enlarged to a 4-input multiplexer to add data sources for the register file: the value provided by the instructions an immediate value, {\tt v}, and the program counter, {\tt pc}, stored in register file as the return address from subroutine.

\includeverilogcode[caption={Register with ALU RALUtoyRISC. File: RALUtoyRISC.sv.  (SystemVerilog)}, 
                    label={lst:RALUtoyRISC}, 
                    firstline=1, 
                    firstnumber=1, 
                    lastline=68]
                    {\verilogSourceDir/A17_RALUtoyRISC.sv}     

\paragraph{Register File} defined in \ref{registerFile}, with the generic design in Listing \ref{lst:A09_regFile.sv}  is adapted in Listing \ref{lst:registerFile} to take into account the interrupt signal, {\tt int} when it is acknowledged by {\tt inta}.

\includeverilogcode[caption={registerFile. File: registerFile.sv.  (SystemVerilog)}, 
                    label={lst:registerFile}, 
                    firstline=1, 
                    firstnumber=1, 
                    lastline=68]
                    {\verilogSourceDir/A17_registerFile.sv}      

\paragraph{ALU} defined in \ref{alu1} and with a generic version described in Listing \ref{lst:A08_alu.sv} is adapted for the instruction set of toyRISC in Listing \ref{lst:ALU}.

\includeverilogcode[caption={ALU. File: ALU.sv.  (SystemVerilog)}, 
                    label={lst:ALU}, 
                    firstline=1, 
                    firstnumber=1, 
                    lastline=68]
                    {\verilogSourceDir/A17_ALU.sv} 

\paragraph{Big Multiplexer} defined \ref{genericMux} is adapted in Listing \ref{lst:bigMux} to take into account the interrupt signal, {\tt int} when it is acknowledged by {\tt inta}.

\includeverilogcode[caption={bigMux. File: bigMux.sv.  (SystemVerilog)}, 
                    label={lst:bigMux}, 
                    firstline=1, 
                    firstnumber=1, 
                    lastline=68]
                    {\verilogSourceDir/A17_bigMux.sv} 

\placedrawing[H]{figures/A17_RISCdetailed.lp}{Top module of toyRISC processor.}{toyRiscProc}

Figure \ref{toyRiscProc} contains the detailed schematic of toyRISC as represented by the Vivado tool. This representation is intended to give only an idea about the complexity of the project. In the electronic version of the book, this image can be inspected in more detail using the zoom feature of the editor.


\subsection{Code generator}\label{codeGen}

In assembly language, the toyRISC processor can only be used if there is a binary code generator, however simple, to translate the mnemonic sequences (see Table \ref{}) into executable code. The Listing \ref{} file describes in System Verilog (for the simplicity of the application) the translation mechanism of all the mnemonics defined for touRISC. In addition to the translation into binary itself, the problem of program jumps to labels defined by the program designer is also solved. Since jumps can be defined both at addresses previous to and after the address at which they are invoked, the translation is done in two passes.


\includeverilogcode[caption={RISCcodeGenerator. File: RISCcodeGenerator.sv.  (SystemVerilog)}, 
                    label={lst:RISCcodeGenerator}, 
                    firstline=1, 
                    firstnumber=1, 
                    lastline=450]
                    {\verilogSourceDir/A17_RISCcodeGenerator.sv} 


\subsection{Simulator}\label{riscsim}

The simulator used to verify the operation of the toyRISC processor, described in the Listing file \ref{lst:testRISC}, generates the context in which the processor works by accessing a program memory and a data memory. The simulator uses the code generator to load the processor's program memory with the executable code. Then the execution of the loaded program is triggered by resetting the processor. If the program is correctly written and the processor has no design errors, then the simulation stops after a finite number of steps with the processor executing the {\tt HALT} instruction. The simulator displays the trace left in the processor structure by the execution of the program by monitoring the program counter, the content of some registers in the register file, and the evolution of the interrupt system signals. The user can add additional signal monitoring to the simulator.


\includeverilogcode[caption={testRISC. File: testRISC.sv.  (SystemVerilog)}, 
                    label={lst:testRISC}, 
                    firstline=1, 
                    firstnumber=1, 
                    lastline=450]
                    {\verilogSourceDir/A17_testRISC.sv} 



\subsection{Testing}

In the file Listing \ref{lst:program} the user can write programs in assembly language. The executable code generation program calls the contents of this file twice (see lines 270 and 272 of {\tt RISCcodeGenerator.sv} in Listing \ref{lst:RISCcodeGenerator}).

\includeverilogcode[caption={Programs to validate the main types of instructions. File: program.sv.  (SystemVerilog)}, 
                    label={lst:program}, 
                    firstline=1, 
                    firstnumber=1, 
                    lastline=450]
                    {\verilogSourceDir/A17_program.sv} 

\begin{exemplu}
Listing \ref{lst:program}  illustrate by its active code the way the interrupt signal act on the toyRISC processor. The simulator activate from the beginning the interrupt signal (see line 17 in Listing\ref{lst:testRISC}). It si initially ignored because the reset signal send  the interrupt enable on {\tt di}. When, in Listing \ref{lst:testRISC} at step 30 the instruction EI (enable interrupt) is executed, the interrupt acts making the program to jump at the instruction stored in the program memory at the location 31 in register file, location loaded with 10 in the first line of the program. In the same time the current PC is saved in the location 30 of the register file. The first instruction in the subroutine located at {|tt PC = 10} disables the interrupt, the next instruction  loads 44 in register 3, and then the instruction {\tt RET} uses the address store in {\tt RF[30]} to return the execution of the program where it was interrupted. The program halts after few steps.
{\small
\begin{verbatim}
progMemory[0]    = 01011111111000000000000000001010
progMemory[1]    = 01011100010000000000000000010111
progMemory[2]    = 01011100000000000000000000001101
progMemory[3]    = 00100000000000000000000000000000
progMemory[4]    = 11001000000000000000000000000010
progMemory[5]    = 11001000000000000000000000000000
progMemory[6]    = 11001000000000000000000000000100
progMemory[7]    = 00011000000000000000000000000000
progMemory[8]    = 11001000000000000000000000000000
progMemory[9]    = 11001000000000000000000000000000
progMemory[10] 	 = 00100100000000000000000000000000
progMemory[11] 	 = 01011100011000000000000000101100
progMemory[12] 	 = 00010100000111100000000000000000
progMemory[13] 	 = xxxxxxxxxxxxxxxxxxxxxxxxxxxxxxxx

pc=  0  RF=[ x,x, x, x, ... x,10] ei=0 inta=0
pc=  1  RF=[ x,x, x, x, ... x,10] ei=0 inta=0
pc=  2  RF=[ x,x,23, x, ... x,10] ei=0 inta=0
pc=  3  RF=[13,x,23, x, ... x,10] ei=0 inta=0
pc=  4  RF=[13,x,23, x, ... x,10] ei=1 inta=1
pc= 10  RF=[13,x,23, x, ... 4,10] ei=0 inta=0
pc= 11  RF=[13,x,23, x, ... 4,10] ei=0 inta=0
pc= 12  RF=[13,x,23,44, ... 4,10] ei=0 inta=0
pc=  4  RF=[13,x,23,44, ... 4,10] ei=0 inta=0
pc=  5  RF=[15,x,23,44, ... 4,10] ei=0 inta=0
pc=  6  RF=[15,x,23,44, ... 4,10] ei=0 inta=0
pc=  7  RF=[19,x,23,44, ... 4,10] ei=0 inta=0
\end{verbatim}
} 

$\diamond$
\end{exemplu}

\subsection{How to proceed}